\documentclass[11pt]{article}
\usepackage[T1]{fontenc}
\usepackage{amsmath, amssymb}
\usepackage{graphicx}
\usepackage{booktabs}
\usepackage{float}
\usepackage[numbers,sort&compress]{natbib}
\usepackage[hidelinks]{hyperref}

\title{Title of the Paper}
\author{Author Name}
\date{\today}

\begin{document}
\maketitle

\begin{abstract}
State the problem, the approach, the main results and what they mean.
\end{abstract}

\section{Introduction}

Give the background and the question. Cite a work as \cite{lecun2015deep}
and several as \cite{lecun2015deep,knuth1984texbook}.

\section{Methods}

\begin{equation}
  y = f(\mathbf{x}; \theta) + \varepsilon
  \label{eq:model}
\end{equation}
In Equation~\eqref{eq:model}, $\theta$ are the parameters.

\section{Results}

\begin{table}[H]
  \centering
  \caption{Results}
  \label{tab:results}
  \begin{tabular}{lcc}
    \toprule
    Method & Metric A & Metric B \\
    \midrule
    Method 1 & 0.91 & 0.88 \\
    Method 2 & 0.94 & 0.90 \\
    \bottomrule
  \end{tabular}
\end{table}

Table~\ref{tab:results} lists the results.

\section{Conclusion}

Summarize the findings and the next steps.

\bibliographystyle{plainnat}
\bibliography{refs}

\end{document}
